\begin{table}[htbp]
\centering
\caption{Price equivalent of moving one position up}
\label{tab:price_equivalent}
\begin{tabular}{lccc}
\toprule
 & (1) Purchases only & \multicolumn{2}{c}{(2) Purchases and clicks} \\
\cmidrule(lr){3-4}
Move & Total & Taste & Search cost \\
\midrule
2 $\to$ 1 & 106.46 & 26.83 & 65.90 \\
3 $\to$ 2 & 54.41 & 15.07 & 35.22 \\
4 $\to$ 3 & 36.59 & 10.52 & 24.09 \\
5 $\to$ 4 & 27.56 & 8.09 & 18.32 \\
10 $\to$ 9 & 12.35 & 3.76 & 8.35 \\
15 $\to$ 14 & 7.96 & 2.45 & 5.41 \\
20 $\to$ 19 & 5.87 & 1.82 & 4.00 \\
25 $\to$ 24 & 4.65 & 1.44 & 3.17 \\
\bottomrule
\multicolumn{4}{p{0.9\linewidth}}{\footnotesize \textit{Note:} The price increase in US dollars that keeps the booking probability constant when a hotel moves one position up, $\Delta p = \tilde p\,(e^{\Delta \log p} - 1)$ with $\Delta \log p = -\frac{\hat\theta_{pos}}{\hat\theta_{price}}\left[\log(r-1) - \log r\right]$, evaluated at the median price, $\tilde p = 129$ USD. Column (1) uses $\hat\theta$ from the logit on bookings. Column (2) splits it into the taste part, $\hat\beta_{pos}$, and the search cost part, $-\hat\gamma_{pos}$. The two parts add up to column (1) in log price, not exactly in dollars.} \\
\end{tabular}
\end{table}
